\section*{\S~2d Funktionen im erweiterten Stab} 
\label{sec:funktionen_im_erweiterten_stab}
\addcontentsline{toc}{section}{\S~2d Funktionen im erweiterten Stab}

\begin{enumerate}[label=(\arabic*)]
    \item Die Unteroffiziere mit Portepee erfüllen im Stab besondere Aufgaben. Diese sind im Nachfolgenden erklärt. Jeder Feldwebel muss eine Funktion besetzen, aber nicht jede Funktion muss besetzt sein. Jeder Feldwebel kann maximal eine Funktion gleichzeitig ausüben.
    \begin{enumerate}[label=(\Alph*)]
        \item \textbf{Der S1-Feldwebel (Innere Führung)} ist Vorsitzender des erweiterten Stabes. Seine Aufgabe besteht darin, die Leitung, den Stab und sonstige zu seiner Amtszeit bestehende Gremien zu kontrollieren. Er leitet außerdem den erweiterten Stab. Sind Unit-Leiter und Spieß abwesend und/oder die Posten nicht besetzt, so übernimmt der S1-Feldwebel bis der Unit-Leiter und der Spieß anwesend und/oder die Posten wieder besetzt sind die Leitung der Unit. Er dient außerdem als oberste Instanz zur Einreichung von Beschwerden gegen die Leitung. Er ist kein Mitglied im Stab.
        \item \textbf{Der Personal-Feldwebel (PersFw)} unterstützt den Gefreiten vom Dienst (GvD) in seiner Funktion als Repräsentant der Mannschaft im erweiterten Stab. Außerdem kontrolliert er den Gefreiten vom Dienst (GvD). Während seiner Amtszeit verzichtet er auf sein Recht, Beschwerde gegen Einzelpersonen, die einen Mannschaftsdienstgrad inne haben, einzureichen. Außerdem unterstützt er die Leitung bei der Mitgliederverwaltung. Er ist kein Mitglied im Stab.
        \item \textbf{Der S6-Feldwebel (Technik)} unterstützt die Leitung in technischen Fragen. Er verwaltet sämtliche öffentlichen Kanäle der 1st Aircat Division administrativ mit. Außerdem unterstützt er Mitglieder bei technischen Problemen im Zusammenhang mit Arma 3 oder den im Modpack festgelegten Mods. Er ist kein Mitglied im Stab.
    \end{enumerate}
    \item Die Funktionen der Feldwebel werden entsprechend der Reihenfolge, wie sie in der Satzung definiert sind, besetzt.
    \item Die Funktionen werden durch Wahl im erweiterten Stab besetzt. Wahlberechtigt sind alle Mitglieder des erweiterten Stabs. Aufgestellt werden können nur Feldwebel. Die Wahl findet jeweils nach der ersten Mission in einem Quartal statt. Es können beliebig viele Amtszeiten, aber nicht mehr als vier, aufeinander folgen, sofern sich dies durch die Anzahl der Feldwebel verhindern lässt. Diese Beschränkung gilt nicht für den S6-Feldwebel, da hier eine besondere technische Expertise notwendig ist. Für die Durchführung von Wahlen gilt \hyperref[sec:abstimmungen_und_wahlen]{\S~3}. Die Wahlleitung übernimmt der Unteroffizier vom Dienst (UvD).
\end{enumerate}